\documentclass[11pt]{article}

% ===================== PACKAGES =====================
\usepackage[a4paper,margin=1in]{geometry}
\usepackage{amsmath,amssymb}
\usepackage{enumitem}
\usepackage{fancyhdr}
\usepackage{xcolor} 

% ===================== HEADER & FOOTER =====================
\pagestyle{fancy}
\fancyhf{}
\lhead{CSE 400: Fundamentals of Probability in Computing}
\rhead{Lecture 19 Scribe}
\cfoot{\thepage}

\title{\textbf{CSE400 -- Fundamentals of Probability in Computing}\\
\large Lecture 19: N-Variate Joint PDF/PMF/CDF}
\author{\textbf{Name:} Maitri Lad
\textbf{Enrollnment no.:} AU2440146}
\date{March 17, 2026}

\begin{document}

\maketitle

\section{Joint Distribution of Random Variables}

\subsection{Definition}

For random variables $X_1, X_2, \dots, X_N$, define:

\[
\mathbf{X} =
\begin{bmatrix}
X_1 \\
X_2 \\
\vdots \\
X_N
\end{bmatrix}
\]

The joint distribution describes the probability behavior of all variables together.

It can be represented using:

\[
\text{Joint PMF, Joint PDF, Joint CDF}
\]

\section{Joint Gaussian Random Variables}

\subsection{Joint Gaussian PDF}

For two variables $X$ and $Y$:

\[
f_{X,Y}(x,y)
=
\frac{1}{2\pi \sigma_X \sigma_Y \sqrt{1-\rho^2}}
\exp \left(
-\frac{1}{2(1-\rho^2)}
\left[
\left(\frac{x-\mu_X}{\sigma_X}\right)^2
+
\left(\frac{y-\mu_Y}{\sigma_Y}\right)^2
-
2\rho
\left(\frac{x-\mu_X}{\sigma_X}\right)
\left(\frac{y-\mu_Y}{\sigma_Y}\right)
\right]
\right)
\]

\subsection{Parameters}

\begin{itemize}
\item $\mu_X, \mu_Y$ : means
\item $\sigma_X, \sigma_Y$ : standard deviations
\item $\rho$ : correlation coefficient
\end{itemize}

\section{Example 1: Joint Gaussian Variables}

\subsection{Given}

\[
\mu_X = 2, \quad \sigma_X^2 = 1
\]

\[
\mu_Y = -3, \quad \sigma_Y^2 = 4
\]

\[
\text{cov}(X,Y) = -1
\]

\subsection{Goal}

\begin{itemize}
\item Write joint PDF $f_{XY}(x,y)$
\item Find marginal PDFs $f_X(x), f_Y(y)$
\item Find $P(X < 0)$
\end{itemize}

\section{Solution (a): Joint PDF Construction}

\subsection{Step 1: Correlation Coefficient}

\[
\rho_{XY}
=
\frac{\text{cov}(X,Y)}{\sigma_X \sigma_Y}
\]

From given variances:

\[
\sigma_X = \sqrt{1} = 1, \quad \sigma_Y = \sqrt{4} = 2
\]

Substitute:

\[
\rho_{XY}
=
\frac{-1}{(1)(2)}
=
-\frac{1}{2}
\]

\subsection{Step 2: Substitute into General Formula}

Using:

\[
\mu_X = 2, \quad \sigma_X = 1
\]

\[
\mu_Y = -3, \quad \sigma_Y = 2
\]

\[
\rho = -0.5
\]

\subsection{Final Joint PDF}

\[
f_{XY}(x,y)
=
\frac{1}{2\pi \sqrt{3}}
\exp \left\{
-\frac{2}{3}
\left[
(x-2)^2
+
\frac{(x-2)(y+3)}{2}
+
\frac{(y+3)^2}{4}
\right]
\right\}
\]

\section{Solution (b): Marginal PDFs}

\subsection{Definition}

\[
f_X(x)
=
\int_{-\infty}^{\infty}
f_{XY}(x,y)\,dy
\]

\[
f_Y(y)
=
\int_{-\infty}^{\infty}
f_{XY}(x,y)\,dx
\]

\subsection{Key Property}

Since $X$ and $Y$ are jointly Gaussian:

\[
X \sim \mathcal{N}(\mu_X, \sigma_X^2)
\]

\[
Y \sim \mathcal{N}(\mu_Y, \sigma_Y^2)
\]

\subsection{Final Marginals}

\[
f_X(x)
=
\frac{1}{\sqrt{2\pi}}
\exp\left(-\frac{(x-2)^2}{2}\right)
\]

\[
f_Y(y)
=
\frac{1}{2\sqrt{2\pi}}
\exp\left(-\frac{(y+3)^2}{8}\right)
\]

\section{Solution (c): CDF and Standardization}

\subsection{Step 1: Identify the CDF}

\[
F_X(x) = P(X \le x)
\]

\[
P(X < 0) = F_X(0)
\]

\subsection{Step 2: Standardization}

\[
X \sim \mathcal{N}(2,1)
\]

\[
Z = \frac{X - \mu_X}{\sigma_X}
\]

\[
P(X < 0)
=
P\left(\frac{X-2}{1} < \frac{0-2}{1}\right)
=
P(Z < -2)
=
\Phi(-2)
\]

\section{Solution (c): Symmetry and Q-function}

\subsection{Step 3: Symmetry Property}

\[
\Phi(-x) = 1 - \Phi(x) = P(Z > x)
\]

\[
P(Z < -2) = P(Z > 2)
\]

\subsection{Step 4: Q-function Relation}

\[
Q(x) = P(Z > x)
\]

\subsection{Final Answer}

\[
P(X < 0) = Q(2)
\]

\section{Application: Large MIMO System}

\subsection{System Setup}

Consider a wireless communication system with:

\[
M \text{ transmit antennas}, \quad N \text{ receive antennas}
\]

\section{Summary}

\subsection{Joint Gaussian}

\[
f_{X,Y}(x,y)
=
\frac{1}{2\pi \sigma_X \sigma_Y \sqrt{1-\rho^2}}
\exp(\cdots)
\]

\subsection{Marginals}

\[
X \sim \mathcal{N}(\mu_X, \sigma_X^2), \quad
Y \sim \mathcal{N}(\mu_Y, \sigma_Y^2)
\]

\subsection{Probability Result}

\[
P(X < 0) = Q(2)
\]

\end{document}